\documentclass[12pt,a4paper]{article}
\usepackage[utf8]{inputenc}
\usepackage[german]{babel}
\usepackage[T1]{fontenc}
\usepackage{newtxtext,newtxmath}
\usepackage{microtype}
\usepackage{enumitem}
\usepackage{longtable,array}
\usepackage{listings,xcolor}
\usepackage[hidelinks]{hyperref}
\usepackage{xurl}
\hypersetup{pdftitle={Installationshandbuch tails-pdp},pdfauthor={Christian Koch}}

\usepackage[left=4cm,right=2cm,top=3cm,bottom=3cm,footskip=1.75cm]{geometry}

% FTK-Layout: einzeiliger Text mit einer Leerzeile zwischen Absätzen.
\usepackage[skip=\baselineskip,indent=0pt]{parskip}
\usepackage{titlesec}
\titleformat{\section}{\normalfont\raggedright\fontsize{16}{19.2}\selectfont\bfseries}{\thesection}{1em}{}
\titleformat{\subsection}{\normalfont\raggedright\fontsize{14}{16.8}\selectfont\bfseries}{\thesubsection}{1em}{}
\titleformat{\subsubsection}{\normalfont\raggedright\fontsize{12}{14.4}\selectfont\bfseries}{\thesubsubsection}{1em}{}
\titleformat{\paragraph}[block]{\normalfont\raggedright\fontsize{12}{14.4}\selectfont\itshape}{\theparagraph}{1em}{}
\titlespacing*{\section}{0pt}{0pt}{\baselineskip}
\titlespacing*{\subsection}{0pt}{2\baselineskip}{\baselineskip}
\titlespacing*{\subsubsection}{0pt}{2\baselineskip}{\baselineskip}
\titlespacing*{\paragraph}{0pt}{2\baselineskip}{\baselineskip}

% Tabellen und Beschriftungen einheitlich in 10 pt.
\usepackage{etoolbox}
\AtBeginEnvironment{tabular}{\fontsize{10}{12}\selectfont}
\usepackage{caption}
\DeclareCaptionFont{ftk}{\fontsize{10}{12}\selectfont}
\captionsetup{font=ftk,format=hang,justification=raggedright,singlelinecheck=false}
\captionsetup[table]{position=top}
\captionsetup[figure]{position=bottom}

% Seitenzahl rechts unten; Grundschrift der Fußzeile 10 pt, Seitenzahl 12 pt.
\usepackage{fancyhdr}
\fancypagestyle{ftk}{
    \fancyhf{}
    \fancyfoot[R]{\normalfont\fontsize{10}{12}\selectfont{\fontsize{12}{14.4}\selectfont\thepage}}
    \renewcommand{\headrulewidth}{0pt}
    \renewcommand{\footrulewidth}{0.4pt}
}
\pagestyle{ftk}
\setlength{\emergencystretch}{3em}

% Codebeispiele in 10 pt, passend zur Tabellenschrift der Thesis.
\AtBeginEnvironment{longtable}{\fontsize{10}{12}\selectfont}
\lstset{basicstyle=\ttfamily\fontsize{10}{12}\selectfont,
  breaklines=true,columns=fullflexible,keepspaces=true,
  showstringspaces=false,frame=single,framerule=0.4pt,
  backgroundcolor=\color{white},aboveskip=\baselineskip,belowskip=\baselineskip}

\title{Installationshandbuch\\tails-pdp\\\large ASBAC-Prototyp im Linux-Kernel}
\author{Christian Koch}
\date{Stand: 25. Oktober 2026}
\begin{document}
\hypersetup{pageanchor=false}
\pagenumbering{gobble}
\maketitle
\thispagestyle{empty}

\clearpage
\pagenumbering{roman}
\hypersetup{pageanchor=true}
\tableofcontents

\clearpage
\pagenumbering{arabic}

\section{Zweck und Geltungsbereich}
Dieses Handbuch beschreibt den Build und die Inbetriebnahme des Prototyps auf einem
vorbereiteten Linux-Zielsystem in drei Arbeitsschritten. Es bezieht sich auf die
bereinigte Repository-Struktur mit den Einstiegsskripten unter \path{scripts/}.
Der Prototyp kontrolliert neue Dateiöffnungen über den BPF-LSM-Hook
\texttt{file\_open} und bewertet bereits offene Dateideskriptoren im Userspace nach.
Socket-Bind-Policies werden nicht unterstützt.

Das Zielsystem wird vor der Installation bereitgestellt. Die Einrichtung des
Hypervisors und die Installation des Betriebssystems sind Voraussetzungen dieses
Ablaufs. Die Referenzkonfiguration beschreibt eine dedizierte Test-VM; die
Installation ist für eine solche Umgebung vorgesehen. Die nachfolgenden Befehle
werden in Bash ausgeführt. Platzhalter wie \texttt{REPOSITORY-URL} und
\texttt{ABGABE-TAG} sind durch die tatsächlich bereitgestellten Werte zu ersetzen.

\section{Voraussetzungen und Lieferumfang}
\subsection{Linux-Zielsystem}
Als Referenz dienen NixOS auf x86\_64 und Linux 6.16.12. Die in der Evaluation
verwendete VM besitzt vier virtuelle CPUs, 8192 MiB RAM und eine 64-GiB-Systemplatte.
Diese Ausstattung ist eine Referenz, keine gemessene Mindestanforderung.
Andere Kernel und Architekturen benötigen eine eigene Funktionsprüfung.

Benötigt werden:
\begin{itemize}[nosep]
  \item ein Kernel mit eBPF, BPF-LSM und Kernel-BTF;
  \item \texttt{bpf} in der aktiven LSM-Reihenfolge;
  \item lesbare BTF-Daten unter \path{/sys/kernel/btf/vmlinux};
  \item eingehängte Dateisysteme \texttt{securityfs} unter
    \path{/sys/kernel/security} und \texttt{bpf} unter \path{/sys/fs/bpf};
  \item Root-Zugriff über \texttt{sudo} für Laden, Attach, Map-Zugriff und
    \texttt{ptrace}-basierten FD-Entzug;
  \item Git, Nix mit \texttt{nix-shell}, Bash und die lokal vorbereiteten
    Buildwerkzeuge; für die geführte Demo außerdem \texttt{tmux}, Python 3,
    \texttt{flock}, \texttt{unshare} und \texttt{mount}.
    Ein Netzwerkzugang kann zur vorgelagerten Beschaffung der Werkzeuge nötig sein;
    der Cargo-Build verwendet die mitgelieferten Abhängigkeitsquellen.
\end{itemize}

Die Dateien \path{nix-config/configuration.nix} und
\path{nix-config/hardware-configuration.nix} dokumentieren die Referenzinstallation.
Die Hardwaredatei gehört zum konkreten Gast und darf auf einem anderen Rechner
nicht ungeprüft dessen Hardwarekonfiguration ersetzen. Die enthaltenen
Initialpasswörter \texttt{password} sind Demowerte; Zugangsdaten des bereitgestellten
Zielsystems werden separat eingerichtet. Die benötigten Kerneloptionen und
LSM-Einstellungen sind durch die Systemadministration vorab bereitzustellen.

Die aktive Umgebung wird geprüft mit:
\begin{lstlisting}
uname -m
uname -r
cat /sys/kernel/security/lsm
test -r /sys/kernel/btf/vmlinux
findmnt -t bpf /sys/fs/bpf
sudo -v
\end{lstlisting}
In der LSM-Ausgabe muss \texttt{bpf} als Eintrag vorhanden sein. Im ausgewerteten
Referenzlauf lautete die Reihenfolge \texttt{capability,landlock,yama,bpf}.
Fehlende Kernelvoraussetzungen werden vor dem Build behoben. Falls nur das
BPF-Dateisystem noch nicht eingehängt ist, kann die Administration es mit
\texttt{sudo mount -t bpf bpf /sys/fs/bpf} einhängen.

\subsection{Buildumgebung und Dateien}
Die Entwicklungsumgebung wird über \path{nix-config/shell.nix} bereitgestellt:
LLVM 21 einschließlich Laufzeitbibliothek, Rustup, Bash, Coreutils, GNU Grep,
Procps, Python 3 und tmux. Die Dateiinhalte von \path{nix-config/rust-toolchain}
und \path{nix-config/ebpf-toolchain} legen die Toolchains fest: aktuell
\texttt{1.93.0} für Userspace und \texttt{nightly-2026-01-15} für eBPF.
Benötigt werden außerdem Rustfmt und Clippy für die Userspace-Toolchain,
\texttt{rust-src} für die eBPF-Toolchain sowie \texttt{bpf-linker} 0.9.15
mit LLVM-21-Unterstützung. Der aktuelle Aya-eBPF-Build benötigt weiterhin Nightly.
Die Nix-Shell setzt \texttt{RUSTUP\_AUTO\_INSTALL=0} und installiert weder Toolchains
noch den Linker beim Start nach.

\path{Cargo.lock} legt die Projektabhängigkeiten fest. Die Quellen unter
\path{vendor/} umfassen Registry- und Git-Abhängigkeiten einschließlich Aya
sowie die Registry-Pakete für den eBPF-Systembibliotheksbuild. Dessen Liste liegt
unter \path{nix-config/ebpf-sysroot/}. Die Quellenersetzung in
\path{.cargo/config.toml} ist bereits aktiviert; Cargo arbeitet standardmäßig offline.
\path{scripts/check-vendor.py} prüft Paketdateien, Prüfsummen und auf Wunsch
auch die Abhängigkeitsliste der installierten \texttt{rust-src}-Komponente.

Die Shell importiert \texttt{<nixpkgs>} ohne feste Revision. Compiler, Rust-Quellen,
Linker, LLVM-Bibliotheken und benötigte Nix-Store-Pfade müssen vorher lokal vorhanden
sein. Ein vollständiges Offline-Werkzeugpaket einschließlich Nixpkgs und Nix-Store-
Abhängigkeiten ist derzeit nicht mitgeliefert. Das Cargo-Vendoring allein ersetzt
keine vorbereitete System- und Werkzeugumgebung.

Der Quellumfang umfasst die acht \texttt{tails-pdp*}-Crates, Cargo-Manifeste,
Lockdatei, \path{.cargo/config.toml}, Skripte, Konfigurationsbeispiele und Lizenzen.
\path{tails-pdp-ebpf/src/vmlinux.rs} wird zum Build benötigt.
Das eBPF-Objekt wird durch das Cargo-Buildskript erstellt und in das Hauptprogramm
eingebettet; es muss zur Ausführung nicht als separate Datei installiert werden.
Das Evaluationsarchiv liegt unter \path{eval-tails-pdp/}.

\section{Installation in drei Arbeitsschritten}
\subsection{Schritt 1: Quellstand und Werkzeuge vorbereiten}
Den bereitgestellten Abgabestand beziehen und eine Nix-Shell öffnen:
\begin{lstlisting}
git clone REPOSITORY-URL tails-pdp
cd tails-pdp
git checkout ABGABE-TAG
nix-shell nix-config/shell.nix
\end{lstlisting}
Bei einer Übergabe als Quellarchiv wird dieses stattdessen entpackt und in sein
Wurzelverzeichnis gewechselt. Die Toolchains, der Linker und die Nix-Abhängigkeiten
müssen auf dem Zielsystem bereits bereitgestellt sein. Innerhalb der Nix-Shell prüfen:
\begin{lstlisting}
rustup run "$(cat nix-config/rust-toolchain)" rustc --version
rustup run "$(cat nix-config/ebpf-toolchain)" rustc --version
setup-bpf-tools
bpf-linker -V
\end{lstlisting}
\texttt{setup-bpf-tools} prüft ausschließlich den vorhandenen Linker auf Version
0.9.15; die Funktion installiert nichts. Die reguläre Toolchain wird aus
\path{nix-config/rust-toolchain} gelesen, die eBPF-Toolchain aus
\path{nix-config/ebpf-toolchain}. Ein unversioniertes \texttt{stable} oder
\texttt{nightly} ersetzt diese Festlegung nicht.

Falls die Toolchains bei der vorgelagerten Vorbereitung mit Netzwerkzugang noch
beschafft werden müssen, können die festgelegten Versionen ausdrücklich installiert werden:
\begin{lstlisting}
rustup toolchain install "$(cat nix-config/rust-toolchain)" --component rustfmt --component clippy
rustup toolchain install "$(cat nix-config/ebpf-toolchain)" --component rust-src
\end{lstlisting}
Die passende Linkerinstallation und alle Nix-Pakete sind ebenfalls vor dem
Offline-Betrieb bereitzustellen. Die Shell stellt zusätzlich die Funktion
\texttt{tails-pdp-admintool} bereit. Sie bestimmt die Repository-Wurzel aus der
Nix-Datei und ruft das dort gebaute Admintool mit \texttt{sudo} auf.
Die folgenden Beispiele verwenden den expliziten Binary-Aufruf.

\subsection{Schritt 2: Programme bauen}
Im Repository-Wurzelverzeichnis innerhalb der Nix-Shell ausführen:
\begin{lstlisting}
bash scripts/demo/build.sh
test -x target/release/tails-pdp
test -x target/release/tails-pdp-admintool
\end{lstlisting}
Der gemeinsame Buildhelfer prüft die festgelegten Toolchains, \texttt{rust-src},
die Vendor-Prüfsummen und \texttt{bpf-linker} 0.9.15. Danach baut er Runtime und
Admintool einschließlich des eingebetteten eBPF-Programms mit
\texttt{cargo build --frozen}. \texttt{--frozen} verbindet eine unveränderte
Abhängigkeitsauflösung mit dem Offline-Modus. Fehlende Quellen oder Werkzeuge
führen zum Abbruch und werden nicht automatisch installiert.

Der entsprechende explizite Cargo-Aufruf lautet:
\begin{lstlisting}
CARGO_NET_OFFLINE=true RUSTUP_AUTO_INSTALL=0 cargo build --frozen --release --bin tails-pdp --bin tails-pdp-admintool
\end{lstlisting}
Alternativ ist weiterhin ein ausdrücklicher Offline-Build möglich:
\begin{lstlisting}
CARGO_NET_OFFLINE=true cargo build --locked --offline --release --bin tails-pdp --bin tails-pdp-admintool
\end{lstlisting}
\texttt{CARGO\_NET\_OFFLINE=true} gilt auch für verschachtelte Cargo-Prozesse des
eBPF-Builds. Der Cargo-Direktaufruf ersetzt die zusätzlichen Prüfungen des
Buildhelfers nicht. Toolchains und Systemwerkzeuge gehören nicht zu \path{vendor/}.
Die unter \path{runtime/bin/} bereitgestellten Binaries besitzen einen Nix-spezifischen
dynamischen Interpreter und benötigen dessen Laufzeitabhängigkeiten; sie sind kein
allgemein auf beliebigen Linux-Systemen ausführbares Paket.

\subsection{Schritt 3: Runtime starten und Bereitschaft prüfen}
Vor dem Start sicherstellen, dass keine weitere Runtime aktiv ist:
\begin{lstlisting}
pgrep -x tails-pdp
\end{lstlisting}
Eine vorhandene Instanz zuerst regulär beenden. Bei einem frischen Checkout sind
keine Beispiel-Policies automatisch aktiviert. Verbliebene lokale \texttt{.policy}-Dateien
unter \path{policies/} vorab prüfen. Dann aus der Repository-Wurzel starten:
\begin{lstlisting}
sudo env RUST_LOG=info ./target/release/tails-pdp
\end{lstlisting}
Das Arbeitsverzeichnis ist Teil der Konfiguration: Die Runtime lädt Policies aus
\path{policies/} und Attribute aus \path{attributes/}. Fehlende Verzeichnisse werden
angelegt. Die Runtime bleibt im Vordergrund. Im Log müssen das Beobachten des
Policyverzeichnisses und \texttt{Waiting for Ctrl-C...} erscheinen; Attach-Fehler
dürfen nicht auftreten. Eine zweite Shell kann den Prozess und die aktiven Daten
prüfen:
\begin{lstlisting}
pgrep -x tails-pdp
sudo ./target/release/tails-pdp-admintool show-active
\end{lstlisting}
Ein leeres Policy-Ergebnis ist beim Erststart zulässig. Ein bloß vorhandener Map-Pin
beweist keine laufende Runtime. Mit Strg+C wird die Runtime regulär beendet.
Gepinnte Maps können danach weiterhin vorhanden sein.

Alternativ bündelt \path{scripts/run.sh} den Admintool-Build und den Runtime-Start:
\begin{lstlisting}
./scripts/run.sh
\end{lstlisting}
Das Skript bestimmt selbst die Repository-Wurzel und verwendet über die
Cargo-Konfiguration \texttt{sudo -E} als Runner. Es aktualisiert den Checkout nicht.

\section{Funktionsprobe mit einer temporären Policy}
Die folgende Probe demonstriert die statische Dateikontrolle für \texttt{cat}.
Sie wird bei laufender Runtime in einer zweiten Shell aus der Repository-Wurzel
ausgeführt. Andere aktive Policies können das Ergebnis beeinflussen; für diese
Probe einen sonst leeren Policystand verwenden.
\begin{lstlisting}
printf 'demo\n' > /tmp/tails-pdp-install-demo.txt
cat > /tmp/tails-pdp-install-demo.policy <<'POLICY'
policy "installation demo"
deny
    action == "file_open";
    command == "cat";
    resource.path == "/tmp/tails-pdp-install-demo.txt";
POLICY
bash scripts/publish-config.sh /tmp/tails-pdp-install-demo.policy policies/install-demo.policy
sudo ./target/release/tails-pdp-admintool show-active
cat /tmp/tails-pdp-install-demo.txt
\end{lstlisting}
Die Veröffentlichung wird asynchron validiert und aktiviert. Erst nach
Bestätigung der Aktivierung im Runtime-Log oder Admintool ist die neue Regel
wirksam. \texttt{cat} muss dann mit einer verweigerten Dateiöffnung enden.
Ein unmittelbar nach der Veröffentlichung noch erfolgreicher Aufruf kann vor
der Aktivierung liegen und wird danach wiederholt.

Zum Aufräumen:
\begin{lstlisting}
rm policies/install-demo.policy
\end{lstlisting}
Nach Bestätigung der neuen Policygeneration muss \texttt{cat} wieder erfolgreich
sein. Anschließend:
\begin{lstlisting}
rm /tmp/tails-pdp-install-demo.policy /tmp/tails-pdp-install-demo.txt
\end{lstlisting}
Änderungen an aktiven Policy- und Attributdateien stets über
\path{scripts/publish-config.sh} veröffentlichen, statt sie während des Betriebs
direkt zu überschreiben. Die Beispiel-Policies mit \path{/home/hntr/test.txt}
benötigen eine passende Ressource und gegebenenfalls angepasste UID-/Attributdaten.

\section{Geführte Demonstration}
Auf dem vorbereiteten dedizierten x86\_64-Linux-System kann die vollständige
Terminal-Demo verwendet werden. Eine zuvor manuell gestartete Runtime beenden;
Demo und Evaluationssuiten nicht gleichzeitig ausführen.
\begin{lstlisting}
./scripts/demo.sh
./scripts/demo.sh --auto
./scripts/demo.sh --build
\end{lstlisting}
Ohne Option führt Enter zum nächsten Schritt. \texttt{--auto} führt die Schritte
ohne Eingabepausen aus; \texttt{--build} erzwingt vorher einen neuen Release-Build.
Fehlende Binaries lösen ebenfalls den Build über die Nix-Umgebung aus.
Werkzeuge und Nix-Abhängigkeiten müssen lokal vorliegen; die Demo lädt sie nicht nach.
Ein eingerichtetes System-\texttt{sudo} beziehungsweise Root-Rechte bleiben erforderlich.

Die Demo verwendet tmux und einen privaten Mount-Namespace mit eigenem bpffs.
Vorhandene Host-Pins bleiben erhalten. Die Konfigurationsverzeichnisse und
\path{demo/files/} werden für die Demo vorbereitet; vorhandene Inhalte werden
gesichert und beim regulären Aufräumen wiederhergestellt. Logs liegen unter
\path{demo/runs/}. Einzelheiten und Diagnoseoptionen beschreibt
\path{scripts/demo/README.md}. Ein erfolgreicher Demolauf ist vom archivierten
Evaluationslauf und dessen Messwerten getrennt zu behandeln.

\section{Prüfkette, Fehlerdiagnose und Grenzen}
Die vollständige Prüfung erfolgt ausschließlich auf dem dedizierten Linux-Zielsystem,
nach Beenden der Demonstrationsruntime. Eine weiterhin wirksame Demo-Policy vorher entfernen.
\begin{lstlisting}
./scripts/test.sh
\end{lstlisting}
Die Prüfkette umfasst Formatierung, Veröffentlichungstest, Demo-Syntax- und
Regressionstests, Rust-Tests, Python-Messauswertung, Clippy, Release-Build,
privilegierte Systemtests und eine automatisch geprüfte vollständige tmux-Demo.
Sie erzeugt ein Ergebnisverzeichnis unter \path{test-results/}. Die privilegierten
Tests verwalten temporär bekannte Maps und isolierte Runtime-Fixtures. Vorhandene
Pins und Prozesse können einen Teststart verhindern.

Bei unvereinbaren gepinnten Map-Layouts oder vor einem ausdrücklich vorgesehenen
Testlauf die Runtime zunächst beenden und die bekannten Pins bereinigen:
\begin{lstlisting}
sudo ./scripts/remove-maps.sh
\end{lstlisting}
Das Skript entfernt aktuelle Pins und historische Pins früherer Fassungen.
Es beendet keine Runtime. Es darf deshalb nicht als laufender Konfigurationswechsel
verwendet werden.

\begin{description}[style=nextline]
  \item[Buildfehler oder fehlender Linker]
    Innerhalb der Nix-Shell bauen, \texttt{setup-bpf-tools} ausführen und Rust-Toolchains
    sowie \texttt{rust-src} prüfen. Den vollständigen Compilerfehler auswerten.
  \item[Fehlendes BTF oder fehlgeschlagener LSM-Attach]
    Kernelvoraussetzungen, aktive LSM-Reihenfolge und BTF-Datei prüfen.
    MacOS ist kein Ausführungsziel für die Runtime.
  \item[Fehler beim Pinnen oder Öffnen von Maps]
    BPF-Mount, Root-Berechtigungen und Map-Layout prüfen. Alte Pins nur nach
    Beenden der Runtime entfernen.
  \item[Policy oder Attribut wird nicht übernommen]
    Arbeitsverzeichnis, Dateiendung und Runtime-Log prüfen. Ungültige Updates lassen
    die letzte gültige Generation bestehen; initial ungültige Daten können den
    Start verhindern. Veröffentlichungsbestätigung allein ist keine Aktivierungsbestätigung.
  \item[FD-Entzug oder privilegierter Test schlägt fehl]
    Root-/ptrace-Berechtigungen und den konkreten Testreport prüfen. Die
    Userspace-Nachbewertung ist nicht atomar mit sämtlichen Dateioperationen;
    ein erfolgreicher Einzellauf beweist keine Race-Freiheit.
\end{description}

\section{Dokumentausgabe und Nachweisstand}
Die LaTeX-Datei ist eigenständig und benötigt keine eingebundenen Thesis-Dateien,
Literaturdaten oder Abbildungen. Für den gemeinsamen Build liegen die Quellen unter
\path{../thesis/}; benötigt werden LaTeX, \texttt{latexmk} und Biber.
Aus der Repository-Wurzel ausführen:
\begin{lstlisting}
bash scripts/build-thesis.sh
\end{lstlisting}
Die lokale Laufkonfiguration \texttt{thesis} ruft denselben Helfer auf.
Er baut Thesis und beide Handbücher und kopiert erst nach erfolgreichem Build
aller drei Dokumente die PDFs nach \path{thesis/final/}. Dieses Handbuch erhält dort den
Namen \path{Christian_Koch_Installationshandbuch_tails-pdp.pdf}; die Thesis wird
mit Namen und Arbeitstitel benannt. Die Handbuchquellen bleiben im vorgesehenen
Thesis-/source-Paket. Eine separate Kompilierung dieses Handbuchs ist weiterhin möglich.

Die dokumentierten Programme, Skripte und Pfade wurden mit dem bereinigten
Quellstand abgeglichen. Die vollständige Neuinstallation nach diesem Handbuch
ist noch auf dem vorgesehenen Linux-Zielsystem zu prüfen. Das archivierte
Evaluationsergebnis vom 10. Oktober 2026 zum Commit \texttt{a9fd6c0} belegt den
damaligen Quellstand; spätere Bereinigungen und dieses Handbuch sind damit nicht
automatisch erneut evaluiert. Für die finale Übergabe sind der tatsächliche
Abgabetag und die ausgelieferten Binärartefakte eindeutig zuzuordnen.
\end{document}
